\chapter{Programs — Pow Table}
%%% generated by the blueprint pipeline from TCSlib.Programs.Light.Lib.PowTable
\section{Overview}

Declarations of the Lean module \texttt{TCSlib.Programs.Light.Lib.PowTable}.

\section{Declarations}

\begin{definition}[List of powers]
\lean{Light.powList}
\label{Light.powList}
\leanok
For natural numbers $b$ and $L$, define a list of integers by taking, in order, $b^j$
for each natural number $j$ with $0\le j<L$, viewing each power as an integer.
\end{definition}

\begin{theorem}[Length of the power list]
\lean{Light.length_powList}
\label{Light.length_powList}
\leanok
\uses{Light.powList}
\difficulty{1}
For all natural numbers $b$ and $L$, the list of integers $b^j$ for $j=0,\ldots,L-1$ has
length $L$.
\end{theorem}

\begin{theorem}[Entries of a list of powers]
\lean{Light.getElem_powList}
\label{Light.getElem_powList}
\leanok
\uses{Light.powList}
\difficulty{1}
For natural numbers $b,L,j$, suppose $j$ is less than the length of the list whose
entries, indexed by $k=0,\ldots,L-1$, are the integers $b^k$. The entry at index $j$ of
that list is the integer $b^j$.
\end{theorem}

\begin{definition}[Destination index]
\lean{Light.PowTable.Dest}
\label{Light.PowTable.Dest}
\leanok
Define the destination index to be $0$.
\end{definition}

\begin{definition}[Table length]
\lean{Light.PowTable.Len}
\label{Light.PowTable.Len}
\leanok
\uses{Light.PowTable.Dest}
Define the table length to be $1$.
\end{definition}

\begin{definition}[Table base]
\lean{Light.PowTable.Base}
\label{Light.PowTable.Base}
\leanok
\uses{Light.PowTable.Dest}
Define the table base to be $2$.
\end{definition}

\begin{definition}[Table exponent index]
\lean{Light.PowTable.Expo}
\label{Light.PowTable.Expo}
\leanok
\uses{Light.PowTable.Dest}
Define the exponent index to be $3$.
\end{definition}

\begin{definition}[Table power index]
\lean{Light.PowTable.Power}
\label{Light.PowTable.Power}
\leanok
\uses{Light.PowTable.Dest}
Define the power index to be $4$.
\end{definition}

\begin{definition}[Power table update round]
\lean{Light.powTableRound}
\label{Light.powTableRound}
\leanok
\uses{Light.PowTable.Base, Light.PowTable.Dest, Light.PowTable.Expo, Light.PowTable.Power, Light.Stmt, Light.v}
The update round first stores the value at index $4$ in memory location $0+3$, then sets
the value at index $4$ to its product with the value at index $2$.
\end{definition}

\begin{definition}[Power table loop body]
\lean{Light.powTableBody}
\label{Light.powTableBody}
\leanok
\uses{Light.PowTable.Expo, Light.PowTable.Len, Light.PowTable.Power, Light.Stmt, Light.Stmt.for, Light.k, Light.powTableRound, Light.v}
The power-table body is the statement that first assigns the value $1$ to the power
variable and then executes a for-loop indexed by the exponent, with the length variable
as its bound and the power-table round as its body.
\end{definition}

\begin{definition}[Power table time bound]
\lean{Light.powTableTime}
\label{Light.powTableTime}
\leanok
For each natural number $L$, the power-table time is $17L+8$.
\end{definition}

\begin{definition}[Filled power table state]
\lean{Light.PowTable.Filled}
\label{Light.PowTable.Filled}
\leanok
\uses{Light.SameOutside, Light.State, Light.frame}
For initial memory $\mu$, destination $dst$, table length $L$, base $b$, a list $t$ of
additional local values, and index $j$, the filled power-table condition on a state means
that for some memory $\mu\prime$ the local frame is
$[dst,L,b,j,b^j]$ followed by $t$, every cell $dst+i$ with $i<j$ contains $b^i$, and
$\mu\prime$ agrees with $\mu$ outside the destination interval of length $L$.
\end{definition}

\begin{theorem}[Power table construction]
\lean{Light.powTable_ends}
\label{Light.powTable_ends}
\leanok
\uses{Light.Ends, Light.Ends.for, Light.PowTable.Filled, Light.SameOutside, Light.SameOutside.refl, Light.Seg, Light.Stmt.blockCost, Light.frame, Light.getElem_powList, Light.powList, Light.powTableBody, Light.powTableRound, Light.powTableTime, Light.update_frame_setLocal}
\difficulty{5}
Let a program be run under machine limits with memory size $S$, word bound $W$, and
recursion depth bound, starting at depth $d$. Let $\mu$ be an initial memory, let
$\mathrm{dst},L,b$ be natural numbers, and let $e,p$ be integers and $t$ a list of
integers. If $S\le W$, the segment of length $L$ starting at $\mathrm{dst}$ fits in
memory, and $(b^j)$ is at most $W$ for every $j\le L$, then the power-table procedure
can execute within time $17L+8$ and leave a memory $\mu'$ whose segment at
$\mathrm{dst}$ is $(1,b,b^2,\ldots,b^{L-1})$, with every location outside that segment
unchanged.
\end{theorem}

\begin{theorem}[Power table procedure specification]
\lean{Light.powTable_meets}
\label{Light.powTable_meets}
\leanok
\uses{Light.Meets, Light.Meets.of_body, Light.SameOutside, Light.Seg, Light.frame_append_zeros, Light.powList, Light.powTableBody, Light.powTableTime, Light.powTable_ends}
\difficulty{1}
Let a program $P$ contain the power-table procedure at position $p$, and let the machine
limits have memory size $S$, word bound $W$, and depth bound. For an initial memory
$\mu$ and natural numbers $\mathrm{dst},L,b$, suppose $S\le W$, the segment of length
$L$ starting at $\mathrm{dst}$ fits in memory, and $b^j\le W$ for every $j\le L$. Then
calling the procedure at depth $d$ with inputs $(\mathrm{dst},L,b)$ can finish within
time $17L+8$, producing the powers $1,b,\ldots,b^{L-1}$ in that segment and preserving
all memory locations outside it.
\end{theorem}
